%% ======================================================================
%%  Chapter VI -- Galois extensions of number fields
%% ======================================================================
\clearpage
\folio{86}%
\chapterhead{6}{VI}{Galois extensions of number fields}{GALOIS EXTENSIONS OF NUMBER FIELDS}
\renewcommand{\chaptertitle}{GALOIS EXTENSIONS OF NUMBER FIELDS}

%% ----------------------------------------------------------------------
\sechead{1}{Galois theory}{GALOIS THEORY}

This section is a supplement to the general theory of commutative fields
presented in \hyperref[chap:II]{Chapter II}, \S\S~3, 4, 6, and 7. Given a field $L$ and a set
$G$ of automorphisms of $L$, one sees immediately that the set of all
$x\in L$ such that $\sigma(x)=x$ for all $\sigma\in G$ is a subfield of
$L$, called the \emph{fixed field} of $G$. It is also clear that, for an
extension $L$ of a field $K$, the set of $K$-automorphisms of $L$ is a
group under composition of mappings.

\begin{theorem}
\label{thm:VI.1.1}
Let $L$ be an extension of finite degree $n$ of a field $K$, where $K$ is
finite or of characteristic zero. The following conditions are equivalent:
\begin{alist}
\item $K$ is the fixed field of the group $G$ of $K$-automorphisms of $L$.
\item For every $x\in L$, the minimal polynomial of $x$ over $K$ has all
its roots in $L$.
\item $L$ is generated by the roots of a polynomial with coefficients in
$K$.
\end{alist}
Under the above conditions the group $G$\footnote{The finiteness of $G$ follows from \hyperref[sec:2.4]{Chapter II, §4}, Theorem
\ref{thm:II.4.1}.} of $K$-automorphisms of $L$ is of order $n$.
\end{theorem}

\begin{proof}
To see that (a) implies (b) observe that, for $x\in L$, the polynomial
$\prod_{\sigma\in G}(X-\sigma(x))=P(X)$ is invariant under $G$ (since any
element of $G$ permutes the factors). Therefore, the coefficients of
$P(X)$ belong to $K$. Since $x$ is a root of $P(X)$ ($G$ contains an
identity element), the minimal polynomial of $x$ divides $P(X)$
(\hyperref[sec:2.3]{Chapter II, §3}, \eqref{eq:II.3.5}). Therefore, (a) implies (b).

To prove that (b) implies (c) take a primitive element $x$ of $L$ over
$K$ (\hyperref[sec:2.4]{Chapter II, §4}, corollary of Theorem \ref{thm:II.4.1}). Its minimal
polynomial over $K$ has all its roots in $L$ by (b). Clearly, these roots
generate $L$ over $K$ (since any one of them does).

Finally, let us show that (c) implies (a). By hypothesis $L$ is generated
over $K$ by a finite number of elements $(x^{(1)},\ldots,x^{(q)})$ and by
all their conjugates $(x_{j}^{(i)})$ (\hyperref[sec:2.4]{Chapter II, §4}). It is clear that any
$K$-isomorphism $\sigma$ of $L$ into an extension of $L$ sends each of
these generators to another element of the same set. Therefore,
$\sigma(L)\subset L$. Moreover, by linear algebra, $\sigma(L)=L$, since
$\sigma$ is $K$-linear and injective. In other words, $\sigma$ is a
$K$-automorphism of $L$. It follows from \hyperref[sec:2.4]{Chapter II, §4}, Theorem
\ref{thm:II.4.1} that the group $G$ of $K$-automorphisms of $L$ has
precisely $n$ elements. Let $x\in L$ be invariant under $G$. Then every
$\sigma\in G$ is a $K[x]$-automorphism of $L$. By \hyperref[sec:2.4]{Chapter II, §4}, Theorem
\ref{thm:II.4.1} there are exactly $[L:K[x]]$ $K[x]$-isomorphisms of $L$
into an extension field of $L$. Thus, $n\leq[L:K[x]]$, from which we may
conclude that $n=[L:K[x]]$. Thus $K[x]=K$, and $x\in K$. This proves that
(c) implies (a). That the order of $G$ is $n$ has been proved along the
way.
\end{proof}

\begin{definition}
\label{def:VI.1.1}
If the conditions of Theorem \ref{thm:VI.1.1} are satisfied, $L$ is called
a \emph{Galois extension} of $K$ and $G$ is called the \emph{Galois group}
of $L$ over $K$. If $G$ is abelian (respectively, cyclic), $L$ is called an
\emph{abelian} (respectively, \emph{cyclic}) extension of $K$.
\end{definition}

\begin{statement}{Corollary of Theorem \ref{thm:VI.1.1}}
Let $K$ be a finite field or a field of characteristic zero, let $L$ be an
extension of finite degree $n$ of $K$, and let $H$ be a group of
automorphisms of $L$ such that $K$ is the fixed field of $H$. Then $L$ is
a Galois extension of $K$ and $H$ is the Galois group of $L$ over $K$.
\end{statement}

\begin{proof}
For $x\in L$ the polynomial $\prod_{\sigma\in H}(X-\sigma(x))=P(X)$ is
invariant under $H$. Therefore, $P(X)$ has its coefficients in $K$. Since
$x$ is a root of $P(X)$, the minimal polynomial of $x$ over $K$ divides
$P(X)$. By Theorem \ref{thm:VI.1.1}, (b), $L$ is a Galois extension of
$K$. Let $G$ be the Galois group of $L$ over $K$. We have $H\subset G$ and
$\operatorname{card}(G)=n$ (Theorem \ref{thm:VI.1.1}). Now let $x$ be a
primitive element of $L$ over $K$ (\hyperref[sec:2.4]{Chapter II, §4}, corollary of Theorem
\ref{thm:II.4.1}) and let $P(X)$ be the polynomial constructed above. Since
$n\leq d^{\circ}(P)$ and $d^{\circ}(P)=\operatorname{card}(H)\leq
\operatorname{card}(G)=n$, we see that $H=G$.
\end{proof}

\begin{theorem}
\label{thm:VI.1.2}
Let $K$ be a field which is finite or of characteristic zero, let $L$ be a
Galois extension of $K$, and let $G$ be the Galois group of $L$ over $K$.
To any subgroup $G'$ of $G$ let us associate the fixed field $k(G')$ of
$G'$, and to any subfield $K'$ of $L$ containing $K$ let us associate the
subgroup $g(K')$ of $G$ consisting of all $K'$-automorphisms of $L$.
\begin{alist}
\item The mappings $g$ and $k$ are bijections and are inverses of one
another. They are both decreasing with respect to the inclusion relations
on $G$ and on $L$, i.e.\ they reverse inclusions. The field $L$ is a
Galois extension of any intermediate field $K'$ (i.e.\
$K\subset K'\subset L$).
\item In order that an intermediate field $K'$ be a Galois extension of
$K$, it is necessary and sufficient that $g(K')$ be a normal subgroup of
$G$. In this case the Galois group of $K'$ over $K$ may be identified with
the quotient group $G/g(K')$.
\end{alist}
\end{theorem}

\begin{proof}
For any intermediate field $K'$ and any $x\in L$ the minimal polynomial of
$x$ over $K'$ divides the minimal polynomial of $x$ over $K$. Thus, all
its roots belong to $L$ by Theorem \ref{thm:VI.1.1}, (b), so $L$ is a
Galois extension of $K'$, also by Theorem \ref{thm:VI.1.1}, (b). $K'$ is
the fixed field of the group $g(K')$ of all $K'$-automorphisms of $L$
(Theorem \ref{thm:VI.1.1}, (a)); in other words, $k(g(K'))=K'$. Let $G'$
be a subgroup of $G$. Then $G'$ is the Galois group of $L$ over $k(G')$
(corollary of Theorem \ref{thm:VI.1.1}); this means $G'=g(k(G'))$. The
relations $k(g(K'))=K'$ and $g(k(G'))=G'$ imply that $k$ and $g$ are
bijections and inverses of one another. It is clear that they reverse
inclusions. This proves (a).

Let us prove (b). Let $K'$ be an intermediate field
($K\subset K'\subset L$). For $x\in K'$ the roots of the minimal polynomial
of $x$ over $K$ are the elements of $L$ of the form $\sigma(x)$, for
$\sigma\in G$. According to Theorem \ref{thm:VI.1.1}, (b), in order that
$K'$ be a Galois extension of $K$, it is necessary and sufficient that
$\sigma(x)\in K'$ for all $x\in K'$ and all $\sigma\in G$, i.e.\ that
$\sigma(K')\subset K'$ for all $\sigma\in G$. But, if
$\sigma(K')\subset K'$, if $\tau\in g(K')$, and if $x\in K'$, then
$\sigma^{-1}\tau\sigma(x)=\sigma^{-1}\sigma(x)=x$, so
$\sigma^{-1}\tau\sigma\in g(K')$. In other words, $K'$ Galois over $K$
implies $g(K')$ normal in $G$. Conversely, suppose that $g(K')$ is normal
in $G$. If $x\in K'$, $\sigma\in G$, and $\tau\in g(K')$, then
$\tau\sigma(x)=\sigma\sigma^{-1}\tau\sigma(x)=\sigma(x)$, since
$\sigma^{-1}\tau\sigma\in g(K')$. Thus $\sigma(x)$ is invariant under
every element of $g(K')$, so $\sigma(x)\in K'$. Consequently, $g(K')$
normal in $G$ implies $\sigma(K')\subset K'$ and so $K'$ is Galois over
$K$.

Finally, let us determine the Galois group of $K'$ over $K$. Since
$\sigma(K')\subset K'$ for all $\sigma\in G$ (and even $\sigma(K')=K'$ by
linear algebra), the restriction $\sigma\mid K'$ of $\sigma$ to $K'$ is a
$K$-automorphism of $K'$. Restriction, $\sigma\mapsto\sigma\mid K'$, is a
homomorphism of $G$ to the Galois group $H$ of $K'$ over $K$. Clearly, its
kernel is $g(K')$. Since
$\operatorname{card}(H)=[K':K]=[L:K][L:K']^{-1}
=\operatorname{card}(G)\operatorname{card}(g(K'))^{-1}
=\operatorname{card}(G/g(K'))$, we may conclude that the restriction
homomorphism is surjective. Therefore, $H\cong G/g(K')$.
\end{proof}

\begin{example}\label{exa:VI.1.1}
Quadratic extensions. Let $K$ be a field of characteristic zero, and let $L$ be a quadratic
extension (i.e.\ of degree $2$) of $K$. As in the beginning of \hyperref[sec:2.5]{§5,
Chapter II} one sees that $L$ is of the form $K[x]$, where $x$ is a root of
a polynomial $X^{2}-d$ ($d\in K$ and $d$ non-square). Since the other root
of this polynomial is $-x$, there is a non-trivial $K$-automorphism
$\sigma$ defined by $\sigma(x)=-x$, i.e.\ $\sigma(a+bx)=a-bx$
($a,b\in K$). Clearly $\sigma^{2}=1$ and $K$ is the fixed field for
$\sigma$. Thus $L$ is a Galois extension of $K$ with the cyclic Galois
group $\{1,\sigma\}$ (by Theorem \ref{thm:VI.1.1} and its corollary).
\end{example}

\begin{example}\label{exa:VI.1.2}
Cyclotomic extensions. Let $K$ be a field of characteristic zero, let $z$ be a primitive $n$th
root of unity in an extension of $K$, and let $L=K(z)$. The field $L$ is
called a \emph{cyclotomic extension} of $K$. The minimal polynomial of $z$
over $K$ divides $X^{n}-1$ (\hyperref[sec:2.3]{Chapter II, §3}, \eqref{eq:II.3.5}), so its
roots are $n$th roots of unity and, consequently, powers of $z$
(\hyperref[sec:1.6]{Chapter I, §6}). Thus $L$ is a Galois extension of $K$ by Theorem
\ref{thm:VI.1.1}, (c).
\end{example}

Let $G$ be the Galois group of $L$ over $K$. Any $\sigma\in G$ is
determined by its effect on $z$. More precisely $\sigma(z)$ is a power
$z^{j(\sigma)}$ of $z$ where $j(\sigma)$ is uniquely determined modulo
$n$. For $\sigma,\tau\in G$ one sees that
$\sigma\tau(z)=\sigma(z^{j(\tau)})=\sigma(z)^{j(\tau)}
=z^{j(\sigma)j(\tau)}$, so $j(\sigma\tau)\equiv j(\sigma)j(\tau)$
(modulo $n$). In other words $\sigma\mapsto j(\sigma)$ defines a
homomorphism $G\to(\Z/n\Z)^{*}$. Since $j(\sigma)$ determines $\sigma$,
this homomorphism is injective, and $G$ is abelian. Thus any cyclotomic
extension is abelian. If $n$ is prime, this extension is even cyclic, for
$G$ is isomorphic to a subgroup of
$(\Z/n\Z)^{*}=\F_{n}^{*}$ (\hyperref[sec:1.7]{Chapter I, §7}, Theorem \ref{thm:I.7.1}, (b)).

As any subgroup of an abelian group is normal, any intermediate field $K'$
in a cyclotomic extension is a Galois extension (and an abelian extension)
of $K$ (Theorem \ref{thm:VI.1.2}, (b)). In particular, any subfield of a
cyclotomic field is an abelian extension of $\Q$. Conversely, it can be
shown (the theorem of Kronecker and Weber) that any abelian extension of
$\Q$ is a subfield of a cyclotomic field.

The reader will observe that, with the preceding notations, the
automorphism $\sigma$ raises every $n$th root of unity to its $j(\sigma)$
power, for all the $n$th roots of unity are powers of $z$. Thus
$\sigma\mapsto j(\sigma)$ is independent of the choice of $z$.

\begin{example}\label{exa:VI.1.3}
Finite fields. Let $\F_{q}$ be a finite field ($q=p^{s}$ with $p$ prime). Any extension
of finite degree of $\F_{q}$ is of the form $\F_{q^{n}}$. Its degree is $n$
(\hyperref[sec:1.7]{Chapter I, §7}). The mapping $\sigma:x\mapsto x^{q}$ is an automorphism of
$\F_{q^{n}}$ (\emph{ibid.}, Proposition \ref{prop:I.7.1}) with $\F_{q}$ as
its fixed field (\emph{ibid.}, Theorem \ref{thm:I.7.1}, (c)). For any
$x\in\F_{q^{n}}$ we have $\sigma^{j}(x)=x^{q^{j}}$ and $\sigma^{n}=1$,
since $x\in\F_{q^{n}}$ satisfies the relation $x^{q^{n}}=x$ (\emph{ibid.},
Theorem \ref{thm:I.7.1}, (c)). On the other hand, for $1\leq j\leq n-1$,
$\sigma^{j}\neq1$ since, if $j<n$, there exists $x\in\F_{q^{n}}$ such that
$x^{q^{j}}\neq x$. Thus, $\{1,\sigma,\ldots,\sigma^{n-1}\}$ is a cyclic
group of order $n$. According to the corollary of Theorem
\ref{thm:VI.1.1}, $\F_{q^{n}}$ is a cyclic extension of degree $n$ of
$\F_{q}$. Its Galois group has a canonical generator, the mapping
$x\mapsto x^{q}$. This mapping is called the \emph{Frobenius
automorphism}.
\end{example}

%% ----------------------------------------------------------------------
\sechead{2}{The decomposition and inertia groups}%
        {THE DECOMPOSITION AND INERTIA GROUPS}

In this section $A$ denotes a Dedekind ring, $K$ its field of fractions
(which is assumed of characteristic zero), $K'$ is a Galois extension of
$K$ of degree $n$, $G$ is the Galois group of $K'$ over $K$, and $A'$ is
the integral closure of $A$ in $K'$.

By applying $\sigma\in G$ to an equation of integral dependence (over $A$)
of an element of $A'$, one sees that $\sigma(x)\in A'$. Therefore,
\begin{equation}
A'\text{ is stable under }G,\text{ i.e. }\sigma(A')=A'\text{ for all }
\sigma\in G.
\label{eq:VI.2.1}
\end{equation}
We have shown only that $\sigma(A')\subset A'$. But we also have
$\sigma^{-1}(A')\subset A'$, so $A'=\sigma\sigma^{-1}(A')\subset\sigma(A')$.
From now on we shall leave the details of arguments of this sort to the
reader.

On the other hand, if $\mathfrak{p}$ is a maximal ideal of $A$ and
$\mathfrak{p}'$ a maximal ideal of $A'$ such that
$\mathfrak{p}'\cap A=\mathfrak{p}$ (i.e.\ $\mathfrak{p}'$ appears in the
factorization of $A'\mathfrak{p}$ into a product of prime ideals in $A'$;
cf.\ \hyperref[sec:5.2]{Chapter V, §2}, \hyperref[prop:V.2.1]{Proposition 1}), then, clearly
$\sigma(\mathfrak{p}')\cap A=\mathfrak{p}$, so
$\sigma(\mathfrak{p}')$ also appears in the expression for
$A'\mathfrak{p}$ as a product of prime ideals of $A'$ and with the same
exponent as $\mathfrak{p}'$. We shall say that $\mathfrak{p}'$ and
$\sigma(\mathfrak{p}')$ are \emph{conjugate prime ideals} of $A'$. We are
going to show that all the prime ideals in the prime factorization of
$A'\mathfrak{p}$ in $A'$ are conjugate.

\begin{proposition}
\label{prop:VI.2.1}
If $\mathfrak{p}$ is a maximal ideal of $A$, then the maximal ideals
$\mathfrak{p}_{i}'$ of $A'$ which appear in the expression for
$A'\mathfrak{p}$ as a product of prime ideals in $A'$ (i.e.\ the maximal
ideals $\mathfrak{p}_{i}'$ are characterized by the property
$\mathfrak{p}_{i}'\cap A=\mathfrak{p}$) are all conjugate. They have the
same residual degree $f$ and the same ramification index $e$. Thus
$A'\mathfrak{p}=\bigl(\prod_{i=1}^{g}\mathfrak{p}_{i}'\bigr)^{e}$ and
$n=efg$.
\end{proposition}

\begin{proof}
The assertion concerning the ramification index and the residual degree is
obvious, because an automorphism preserves all algebraic relations. The
formula $n=efg$ is thus a special case of the relation
$\sum e_{i}f_{i}=n$ (\hyperref[sec:5.2]{Chapter V, §2}, \hyperref[thm:V.2.1]{Theorem 1}). Now let $\mathfrak{p}'$ be
one of the $\mathfrak{p}_{i}'$, and assume that another of the
$\mathfrak{p}_{i}'$, which we shall denote $\mathfrak{q}'$, is not
conjugate to $\mathfrak{p}'$. Since $\mathfrak{q}'$ and
$\sigma(\mathfrak{p}')$ ($\sigma\in G$) are maximal and distinct,
$\sigma(\mathfrak{p}')\not\subset\mathfrak{q}'$. We need the following
lemma.
\end{proof}

\begin{lemma}
\label{lem:VI.2.1}
(``lemma for staying away from prime ideals''). Let $R$ be a ring,
$\mathfrak{p}_{1},\ldots,\mathfrak{p}_{q}$ a finite set of prime ideals of
$R$, and let $\mathfrak{b}$ be an ideal of $R$ such that
$\mathfrak{b}\not\subset\mathfrak{p}_{i}$ for any index $i$. Then there
exists $b\in\mathfrak{b}$ such that $b\notin\mathfrak{p}_{i}$ for any $i$.
\end{lemma}

\begin{proof}[Proof of the lemma]
Without loss of generality we may assume
$\mathfrak{p}_{j}\not\subset\mathfrak{p}_{i}$ for $i\neq j$. Let
$x_{ij}\in\mathfrak{p}_{j}-\mathfrak{p}_{i}$ for $i\neq j$,
$1\leq i,j\leq q$. Since $\mathfrak{b}\not\subset\mathfrak{p}_{i}$, there
exists $a_{i}\in\mathfrak{b}-\mathfrak{p}_{i}$. Put
$b_{i}=a_{i}\prod_{j\neq i}x_{ij}$. Then $b_{i}\in\mathfrak{b}$,
$b_{i}\in\mathfrak{p}_{j}$ for $j\neq i$, and
$b_{i}\notin\mathfrak{p}_{i}$ (since $\mathfrak{p}_{i}$ is prime). Thus
\[
b=b_{1}+\cdots+b_{q}\in\mathfrak{b}-\bigcup_{i=1}^{q}\mathfrak{p}_{i},
\]
since clearly $b\in\mathfrak{b}$ and
$b\equiv b_{i}\pmod{\mathfrak{p}_{i}}$ (so $b\notin\mathfrak{p}_{i}$ for
any $i$).
\end{proof}

Returning to the proof of the proposition, we see, from the lemma, that
there is an element $x\in\mathfrak{q}'$ such that
$x\notin\sigma(\mathfrak{p}')$ for all $\sigma\in G$. Consider the norm of
$x$, $\mathrm{N}(x)=\prod_{\tau\in G}\tau(x)$ (\hyperref[sec:2.6]{Chapter II, §6},
Proposition \ref{prop:II.6.1}). Since $\tau(x)\in A'$ for all $\tau\in G$
(by \eqref{eq:VI.2.1}) we see that $\mathrm{N}(x)\in\mathfrak{q}'$; in
fact $\mathrm{N}(x)\in\mathfrak{q}'\cap A=\mathfrak{p}$. On the other
hand $x\notin\tau^{-1}(\mathfrak{p}')$, so
$\tau(x)\notin\mathfrak{p}'$ for any $\tau\in G$. Since
$\mathfrak{p}'$ is prime, we see that
$\mathrm{N}(x)\notin\mathfrak{p}'$, and this contradicts
$\mathrm{N}(x)\in\mathfrak{p}$.

Now let $\mathfrak{p}'$ be a maximal ideal of $A'$ such that
$\mathfrak{p}'\cap A=\mathfrak{p}$. Those $\sigma\in G$ for which
$\sigma(\mathfrak{p}')=\mathfrak{p}'$ form a subgroup $D$ of $G$, called
the \emph{decomposition group} of $\mathfrak{p}'$. If $g$ denotes the
number of conjugates of $\mathfrak{p}'$, then
\begin{equation}
g=\operatorname{card}(G)\operatorname{card}(D)^{-1}
\quad\text{or}\quad
\operatorname{card}(D)=\frac{n}{g}=ef.
\label{eq:VI.2.2}
\end{equation}

For $\sigma\in D$ the relations $\sigma(A')=A'$ and
$\sigma(\mathfrak{p}')=\mathfrak{p}'$ imply that $\sigma$ induces an
automorphism $\bar{\sigma}$ of $A'/\mathfrak{p}'$
($x\equiv y\pmod{\mathfrak{p}'}$ entails
$\sigma(x)\equiv\sigma(y)\pmod{\mathfrak{p}'}$). It is clear that
$\bar{\sigma}$ is an $(A/\mathfrak{p})$-automorphism. The mapping
$\sigma\mapsto\bar{\sigma}$ is a group homomorphism, whose kernel is the
subgroup $I\subset D$ consisting of those $\sigma\in D$ which satisfy the
relation $\sigma(x)-x\in\mathfrak{p}'$ for all $x\in A'$. Thus, $I$ is a
normal subgroup of $D$, called the \emph{inertia subgroup} of
$\mathfrak{p}'$.

\begin{proposition}
\label{prop:VI.2.2}
With the same notations as above, assume that $A/\mathfrak{p}$ is finite or
of characteristic zero. Then $A'/\mathfrak{p}'$ is a Galois extension of
degree $f$ of $A/\mathfrak{p}$, and the mapping
$\sigma\mapsto\bar{\sigma}$ is a surjective homomorphism of $D$ on the
Galois group of $A'/\mathfrak{p}'$ over $A/\mathfrak{p}$. Moreover,
$\operatorname{card}(I)=e$.
\end{proposition}

\begin{proof}
Let $K_{D}$ be the fixed field of $D$, let $A_{D}=A'\cap K_{D}$ be the
integral closure of $A$ in $K_{D}$, and let $\mathfrak{p}_{D}$ be the prime
ideal $\mathfrak{p}'\cap A_{D}$. According to Proposition
\ref{prop:VI.2.1} and the definition of $D$, $\mathfrak{p}'$ is the only
prime factor of $A'\mathfrak{p}_{D}$. Put
$A'\mathfrak{p}_{D}=\mathfrak{p}'^{e'}$ and write $f'$ for the residual
degree $[A'/\mathfrak{p}':A_{D}/\mathfrak{p}_{D}]$. According to
\hyperref[thm:V.2.1]{Theorem 1} of \hyperref[sec:5.2]{§2, Chapter V}, Theorem \ref{thm:VI.1.2} of \hyperref[sec:6.1]{§1}, and
\eqref{eq:VI.2.2}, we have
$e'f'=[K':K_{D}]=\operatorname{card}(D)=ef$. Since
$A/\mathfrak{p}\subset A_{D}/\mathfrak{p}_{D}\subset
A'/\mathfrak{p}'$, $f'\leq f$. Since
$\mathfrak{p}A_{D}\subset\mathfrak{p}_{D}$, $e'\leq e$. Therefore, since
$e'f'=ef$, $e=e'$ and $f=f'$. This entails
\begin{equation}
A/\mathfrak{p}\simeq A_{D}/\mathfrak{p}_{D}.
\label{eq:VI.2.3}
\end{equation}

Now let $\bar{x}$ be a primitive element for $A'/\mathfrak{p}'$ over
$A/\mathfrak{p}$ and let $x\in A'$ be a representative for $\bar{x}$. Let
$X^{r}+a_{r-1}X^{r-1}+\cdots+a_{0}=P(X)$ be the minimal polynomial for
$x$ over $K_{D}$. We know that $a_{i}\in A_{D}$ (\hyperref[sec:2.6]{Chapter II, §6}, the
corollary of Proposition \ref{prop:II.6.2}). The roots of $P(X)$ are all
of the form $\sigma(x)$ with $\sigma\in D$. The ``reduced'' polynomial
$\overline{P}(X)=X^{r}+\bar{a}_{r-1}X^{r-1}+\cdots+\bar{a}_{0}$ has its
coefficients in $A/\mathfrak{p}$ (by \eqref{eq:VI.2.3}) and the roots of
$\overline{P}(X)$ are all of the form $\bar{\sigma}(\bar{x})$ with
$\sigma\in D$. Consequently, $A'/\mathfrak{p}'$ contains all the
conjugates of $\bar{x}$ over $A/\mathfrak{p}$, and
$A'/\mathfrak{p}'$ is a Galois extension of $A/\mathfrak{p}$ (\hyperref[sec:6.1]{§1},
Theorem \ref{thm:VI.1.1}, (c)). Furthermore, since every conjugate of
$\bar{x}$ over $A/\mathfrak{p}$ is of the form
$\bar{\sigma}(\bar{x})$, every $(A/\mathfrak{p})$-automorphism of
$A'/\mathfrak{p}'$ is of the form $\bar{\sigma}$ for some $\sigma\in D$.
Thus, the Galois group of $A'/\mathfrak{p}'$ over $A/\mathfrak{p}$ may be
identified with $D/I$. Since its order is
$[A'/\mathfrak{p}':A/\mathfrak{p}]=f$,
$\operatorname{card}(D)/\operatorname{card}(I)=f$, so
$\operatorname{card}(I)=e$, by \eqref{eq:VI.2.2}.
\end{proof}

\begin{statement}{Corollary}
In order that $\mathfrak{p}$ not ramify in $A'$ it is necessary and
sufficient that the inertia group $I$ of $\mathfrak{p}'$ for any
$\mathfrak{p}'$ over $\mathfrak{p}$ (i.e.\ such that
$\mathfrak{p}'\cap A=\mathfrak{p}$) be trivial.
\end{statement}

\begin{statement}{Remark}
Write $D_{\mathfrak{p}'}$ and $I_{\mathfrak{p}'}$ for the decomposition
and inertia groups of the maximal ideal $\mathfrak{p}'\subset A'$. For a
conjugate ideal $\sigma(\mathfrak{p}')$ ($\sigma\in G$)
\begin{equation}
D_{\sigma(\mathfrak{p}')}=\sigma D_{\mathfrak{p}'}\sigma^{-1}
\quad\text{and}\quad
I_{\sigma(\mathfrak{p}')}=\sigma I_{\mathfrak{p}'}\sigma^{-1}.
\label{eq:VI.2.4}
\end{equation}
To prove \eqref{eq:VI.2.4} note that, for
$\tau\in D_{\mathfrak{p}'}$, we have
\[
\sigma\tau\sigma^{-1}\sigma(\mathfrak{p}')
=\sigma\tau(\mathfrak{p}')=\sigma(\mathfrak{p}'),
\]
so $\sigma D_{\mathfrak{p}'}\sigma^{-1}
\subset D_{\sigma(\mathfrak{p}')}$. The reverse inclusion follows by an
argument analogous to that of the remark following \eqref{eq:VI.2.1}.
Similarly, for $\tau\in I_{\mathfrak{p}'}$ and $x\in A'$,
\[
\sigma\tau\sigma^{-1}(x)-x
=\sigma\tau(\sigma^{-1}(x))-\sigma\sigma^{-1}(x)
=\sigma\bigl(\tau(\sigma^{-1}(x))-\sigma^{-1}(x)\bigr)
\in\sigma(\mathfrak{p}'),
\]
so $\sigma I_{\mathfrak{p}'}\sigma^{-1}
\subset I_{\sigma(\mathfrak{p}')}$, the reverse inclusion being proved as
before.
\end{statement}

When $K'$ is an abelian extension of $K$, the groups
$D_{\sigma(\mathfrak{p}')}$ (respectively,
$I_{\sigma(\mathfrak{p}')}$) ($\sigma\in G$) are all equal. They depend
only on the ideal $\mathfrak{p}$ of the ring $A$ (Proposition
\ref{prop:VI.2.1}).

%% ----------------------------------------------------------------------
\sechead{3}{The number field case. The Frobenius automorphism}%
        {THE NUMBER FIELD CASE. THE FROBENIUS AUTOMORPHISM}

The preceding will now be applied to number fields and their rings of
integers. This application is possible because number fields are of
characteristic zero and their residual fields are finite.

Let us keep our earlier notations ($K\subset K'$, both number fields, $K'$
a Galois extension of $K$ with Galois group $G$, the respective rings of
integers $A$ and $A'$). Let $\mathfrak{p}$ be a maximal ideal of $A$ which
does not ramify in $A'$, and let $\mathfrak{p}'$ be a prime factor of
$A'\mathfrak{p}$. The inertia group of $\mathfrak{p}'$ consists of the
identity of $G$ alone (\hyperref[sec:6.2]{§2}, corollary of Proposition \ref{prop:VI.2.2});
its decomposition group $D$ is canonically isomorphic to the Galois group
of $A'/\mathfrak{p}'$ over $A/\mathfrak{p}$ (\hyperref[sec:6.2]{§2}, Proposition
\ref{prop:VI.2.2}). But the Galois group of $A'/\mathfrak{p}'$ over
$A/\mathfrak{p}$ is cyclic with a canonical generator
$\bar{\sigma}:\bar{x}\mapsto\bar{x}^{q}$, where
$q=\operatorname{card}(A/\mathfrak{p})$ (\hyperref[sec:6.1]{§1}, \hyperref[exa:VI.1.3]{Example 3}). Therefore, $D$
itself is cyclic with a canonical generator $\sigma$ defined by the
relation $\sigma(x)\equiv x^{q}\pmod{\mathfrak{p}'}$ for any $x\in A'$.
This generator is called the \emph{Frobenius automorphism} of
$\mathfrak{p}'$. We denote it $(\mathfrak{p}',K'/K)$.

For $\tau\in G$ we have, as in the remark at the end of \hyperref[sec:6.2]{§2},
\begin{equation}
(\tau(\mathfrak{p}'),K'/K)
=\tau(\mathfrak{p}',K'/K)\tau^{-1}.
\label{eq:VI.3.1}
\end{equation}
In particular, if $K'$ is an abelian extension,
$(\mathfrak{p}',K'/K)$ depends only on the ideal $\mathfrak{p}$ of $A$.
In this case we write
$\left(\frac{K'/K}{\mathfrak{p}}\right)$.

\begin{proposition}
\label{prop:VI.3.1}
With the preceding hypotheses and notations let $F$ be an intermediate
field ($K\subset F\subset K'$) and write $f$ for the residual degree of
$\mathfrak{p}'\cap F$ over $K$. Then
\begin{alist}
\item $(\mathfrak{p}',K'/F)=(\mathfrak{p}',K'/K)^{f}$;
\item if $F$ is Galois over $K$, the restriction of
$(\mathfrak{p}',K'/K)$ to $F$ equals
$(\mathfrak{p}'\cap F,F/K)$.
\end{alist}
\end{proposition}

\begin{proof}
(a) Put $\sigma=(\mathfrak{p}',K'/K)$. By definition
$\sigma(\mathfrak{p}')=\mathfrak{p}'$ and
$\sigma(x)\equiv x^{q}\pmod{\mathfrak{p}'}$ for every $x\in A'$ (here
$q=\operatorname{card}(A/\mathfrak{p})$). Thus,
$\sigma^{f}(\mathfrak{p}')=\mathfrak{p}'$ and
$\sigma^{f}(x)\equiv x^{q^{f}}\pmod{\mathfrak{p}'}$ for all $x\in A'$. By
definition of $f$, $q^{f}$ is the cardinality of the residual field
$(A'\cap F)/(\mathfrak{p}'\cap F)$. Furthermore, the decomposition group
of $\mathfrak{p}'$ over $F$ is obviously a subgroup of the decomposition
group $D$ of $\mathfrak{p}'$ over $K$. It is of order
\[
[A'/\mathfrak{p}':(A'\cap F)/(\mathfrak{p}'\cap F)]
=f^{-1}[A'/\mathfrak{p}':A/\mathfrak{p}]
=f^{-1}\operatorname{card}(D),
\]
by \eqref{eq:VI.2.2} of \hyperref[sec:6.2]{§2}. Since $D$ is cyclic and generated by
$\sigma$, the only subgroup of $D$ of order
$f^{-1}\operatorname{card}(D)$ is generated by $\sigma^{f}$. This proves
(a).

(b) Suppose $F$ is Galois over $K$ and write $\sigma'$ for the
restriction to $F$ of $\sigma$ (\hyperref[sec:6.1]{§1}, Theorem \ref{thm:VI.1.1}, (b)). Since
$\sigma(\mathfrak{p}')=\mathfrak{p}'$, it follows that
$\sigma'(\mathfrak{p}'\cap F)=\mathfrak{p}'\cap F$ and $\sigma'$ belongs
to the decomposition group of $\mathfrak{p}'\cap F$ over $K$. Moreover, it
is clear that
$\sigma'(x)\equiv x^{q}\pmod{\mathfrak{p}'\cap F}$ for all
$x\in A'\cap F$, with $q=\operatorname{card}(A/\mathfrak{p})$.
\end{proof}

%% ----------------------------------------------------------------------
\sechead{4}{An application to cyclotomic fields}%
        {AN APPLICATION TO CYCLOTOMIC FIELDS}

We are going to utilize the theory which we have just developed to present
a third proof of the irreducibility of the cyclotomic polynomial (cf.\
\hyperref[sec:2.9]{Chapter II, §9}, Theorem \ref{thm:II.9.1} and \hyperref[sec:5.2]{Chapter V, §2}, Example).

\begin{theorem}
\label{thm:VI.4.1}
Let $z$ be a primitive complex $n$th root of unity. Then
\begin{alist}
\item No prime number which does not divide $n$ ramifies in $\Q[z]$;
\item $\Q[z]$ is an abelian extension of $\Q$ of degree $\varphi(n)$ and
with Galois group isomorphic to $(\Z/n\Z)^{*}$.
\end{alist}
\end{theorem}

\begin{proof}
(a) Let $F(X)$ be the minimal polynomial of $z$ over $\Q$ and let $d$ be
its degree (so $d=[\Q[z]:\Q]$). The polynomial $F(X)$ divides
$X^{n}-1$; let $X^{n}-1=F(X)G(X)$. We know that
$\mathrm{D}(1,z,\ldots,z^{d-1})=\pm\mathrm{N}(F'(z))$ (\hyperref[sec:2.7]{Chapter II, §7},
\eqref{eq:II.7.6}). From the relation
$nX^{n-1}=F'(X)G(X)+F(X)G'(X)$, we may prove, by substitution,
$nz^{n-1}=F'(z)G(z)$. Since $z$ is a unit of $\Q[z]$, it is of norm
$\pm1$. Taking norms, we may conclude that $\mathrm{N}(F'(z))$ divides
$n^{d}$. Finally, since $z$ is an integer in $\Q[z]$, the absolute
discriminant of $\Q[z]$ divides $\mathrm{D}(1,z,\ldots,z^{d-1})$ and,
therefore also, $n^{d}$. By \hyperref[sec:5.3]{Chapter V, §3}, \hyperref[thm:V.3.1]{Theorem 1}, no prime number
which does not divide $n$ ramifies in $\Q[z]$. This proves (a).

(b) Recall (\hyperref[sec:6.1]{§1}, \hyperref[exa:VI.1.2]{Example 2}) that $\Q[z]$ is an abelian extension of $\Q$
and that there is an injective homomorphism $j$ of the Galois group $G$ of
$\Q[z]$ over $\Q$ into $(\Z/n\Z)^{*}$. More precisely, the element
$\sigma\in G$ raises all the $n$th roots of unity to the power
$j(\sigma)$. Let $p$ be a prime number which does not divide $n$. By (a)
the Frobenius automorphism
$\left(\frac{\Q[z]/\Q}{p}\right)$ is defined; denote it $\sigma_{p}$.
Writing $A$ for the ring of integers of $\Q[z]$ and $\mathfrak{p}$ for an
arbitrary prime factor of $Ap$, we obtain, from the definition of the
Frobenius automorphism, the relation
$\sigma_{p}(x)\equiv x^{p}\pmod{\mathfrak{p}}$ for all $x\in A$. In
particular, putting $j=j(\sigma_{p})$, we have
$z^{f}\equiv z^{p}\pmod{\mathfrak{p}}$. Recall that
\[
\prod_{\substack{0\leq r\leq n-1\\ r\not\equiv p\pmod n}}
(z^{p}-z^{r})=P'(z^{p})=nz^{p(n-1)},
\]
where $P(X)=X^{n}-1=\prod_{0\leq r\leq n-1}(X-z^{r})$. Since $n$ is
relatively prime to $p$, since
$\mathfrak{p}\cap\Z=p\Z$, and since $z$ is a unit in the integers of
$\Q[z]$, we may conclude from the relation
$P'(z^{p})=nz^{p(n-1)}$ that
$\prod_{\substack{0\leq r\leq n-1\\ r\not\equiv p\pmod n}}
(z^{p}-z^{r})\notin\mathfrak{p}$. The relation
$z^{f}\equiv z^{p}\pmod{\mathfrak{p}}$ thus implies that $j$ represents
the residue class of $p\bmod n$. Hence $j(G)$ contains the residue
classes mod $n$ of all the prime numbers $p$ which do not divide $n$. But
this means that $j(G)=(\Z/n\Z)^{*}$, which proves (b).
\end{proof}

%% ----------------------------------------------------------------------
\sechead{5}{Another proof of the quadratic reciprocity law}%
        {ANOTHER PROOF OF THE QUADRATIC RECIPROCITY LAW}

Let $q$ be an odd prime. Let $K$ be the cyclotomic field generated by a
primitive $q$th root of unity in $\C$. The Galois group $G$ of $K$ over
$\Q$ is isomorphic to $\F_{q}^{*}$ (\hyperref[sec:6.4]{§4}, Theorem \ref{thm:VI.4.1}, (b));
it is cyclic of even order $q-1$. There is a unique subgroup $H$ of index
$2$, which corresponds to the subgroup of squares
$(\F_{q}^{*})^{2}\subset\F_{q}^{*}$. Thus $K$ contains a unique quadratic
subfield $F$ (\hyperref[sec:6.1]{§1}, Theorem \ref{thm:VI.1.2}, (b)). No prime number $p\neq
q$ ramifies in $F$; for, if it did, it would ramify in $K$; and this would
contradict Theorem \ref{thm:VI.4.1}, (a) of \hyperref[sec:6.4]{§4}. The calculation of the
discriminant of a quadratic field (\hyperref[sec:5.3]{Chapter V, §3}, \hyperref[exa:V.3.1]{Example 1}) implies that
$F=\Q[\sqrt{q}]$, if $q\equiv1\pmod4$, and $F=\Q[\sqrt{-q}]$, if
$q\equiv3\pmod4$. Putting $q^{*}=(-1)^{(q-1)/2}q$, we have, in either
case, $F=\Q[\sqrt{q^{*}}]$.

Let $p$ be a prime distinct from $q$. Write $\sigma_{p}$ for the
Frobenius automorphism
$\left(\frac{K/\Q}{p}\right)$ (cf.\ \hyperref[sec:6.4]{§4}). Its restriction to $F$ is
$\left(\frac{F/\Q}{p}\right)$ (\hyperref[sec:6.3]{§3}, Proposition \ref{prop:VI.3.1}, (b)).
It is the identity if $\sigma_{p}\in H$, i.e.\ if the exponent
$j(\sigma_{p})=$ the residue class of $p\bmod q$ (cf.\ \hyperref[sec:6.4]{§4}) is a square in
$\F_{q}^{*}$; otherwise it is the non-trivial automorphism of $F$. In
other words, identifying the Galois group $G/H$ of $F$ over $\Q$ with
$\{+1,-1\}$, we have
\begin{equation}
\left(\frac{F/\Q}{p}\right)=\left(\frac{p}{q}\right)
\label{eq:VI.5.1}
\end{equation}
by definition of the Legendre symbol $\left(\frac{p}{q}\right)$
(\hyperref[sec:5.5]{Chapter V, §5}).

On the other hand, the results concerning the decomposition of the prime
$p$ in $F=\Q[\sqrt{q^{*}}]$ (\hyperref[sec:5.4]{Chapter V, §4}) give further information
regarding $\left(\frac{F/\Q}{p}\right)$. By definition it is the identity
if $p$ splits in $F$, and it is the non-trivial automorphism if $p$
remains prime. Using \hyperref[sec:5.4]{Chapter V, §4}, \hyperref[prop:V.4.1]{Proposition 1}, we may conclude that,
if $p$ is odd,
\begin{equation}
\left(\frac{F/\Q}{p}\right)=\left(\frac{q^{*}}{p}\right).
\label{eq:VI.5.2}
\end{equation}

Comparing \eqref{eq:VI.5.1} and \eqref{eq:VI.5.2} we obtain
\[
\left(\frac{p}{q}\right)=\left(\frac{q^{*}}{p}\right)
=\left(\frac{-1}{p}\right)^{(q-1)/2}\left(\frac{q}{p}\right).
\]
But
\[
\left(\frac{-1}{p}\right)=(-1)^{(p-1)/2},
\]
by Euler's criterion (\hyperref[sec:5.5]{Chapter V, §5}, \hyperref[prop:V.5.1]{Proposition 1}). Consequently,
\[
\left(\frac{p}{q}\right)=(-1)^{(p-1)(q-1)/4}\left(\frac{q}{p}\right).
\]
This completes our second proof of the quadratic reciprocity law
(\hyperref[sec:5.5]{Chapter V, §5}, \hyperref[thm:V.5.1]{Theorem 1}).

To take care of the case $p=2$, recall that $2$ splits in $F$ if
$q^{*}\equiv1\pmod8$ and $2$ remains prime if
$q^{*}\equiv5\pmod8$ (\hyperref[sec:5.4]{Chapter V, §4}, \hyperref[prop:V.4.1]{Proposition 1}). However,
$(-1)^{(q^{2}-1)/8}=(-1)^{(q^{*2}-1)/8}$, and this is $1$ if
$q^{*}\equiv1\pmod8$ and $-1$ if $q^{*}\equiv5\pmod8$. Thus,
\begin{equation}
\left(\frac{F/\Q}{2}\right)=(-1)^{(q^{2}-1)/8}.
\label{eq:VI.5.3}
\end{equation}
Putting \eqref{eq:VI.5.1} and \eqref{eq:VI.5.3} together, we obtain
$\left(\frac{2}{q}\right)=(-1)^{(q^{2}-1)/8}$. This is the difficult
``complementary formula'' (\hyperref[sec:5.5]{Chapter V, §5}, \hyperref[prop:V.5.2]{Proposition 2}, (b)).
